\documentclass[10pt,a4paper]{article}
\usepackage[utf8]{inputenc}
\usepackage[T1]{fontenc}
\usepackage{lmodern}
\usepackage[margin=2.2cm]{geometry}
\usepackage{amsmath,amssymb}
\usepackage{graphicx}
\usepackage{booktabs}
\usepackage{caption}
\usepackage{float}
\usepackage{hyperref}
\captionsetup{font=small,labelfont=bf}
\input{values}

\title{How the accuracy of the model is evaluated\\[2pt]
\large from the line map to the charge states}
\author{\small Generated by \texttt{docs/charge\_states/make\_eval.py}\\
\small run: \texttt{\vRun}}
\date{}

\begin{document}
\maketitle

\noindent
This note describes the evaluation, not the results: what is compared with
what, on which devices, with which metric, and against which reference.
The numbers shown are there to illustrate each definition.  The worked
example is test device \vDevice{} with the \vBud{} U-Net.

\section{What is compared with what}

\subsection*{The ground truth}
For every simulated device, QArray computes the lowest-energy occupation
$(n_1,n_2)$ at every pixel of the $100\times100$ gate grid, from the
capacitance model (\texttt{charge\_states.npy}).  No sensor is involved.  The
true line map $Y$ marks the pixels where $(n_1,n_2)$ changes between
neighbours.  Both are known exactly in simulation and never in an
experiment.

\subsection*{The prediction}
The U-Net sees only the charge-sensor signal on the measured rays
(Fig.~\ref{fig:line}a).  It returns a probability $P$ per pixel, and
$B = P > t$ is its line map.  From $B$, \texttt{label\_states} assigns
$(n_1,n_2)$ to every pixel, optionally after \texttt{fit\_lines} has
repaired the gaps.

\subsection*{Two levels of evaluation}
\begin{enumerate}
\item \textbf{Line level}: is $B$ close to $Y$?  This is what the U-Net is
  trained and selected for.
\item \textbf{Charge-state level}: are the labels from $B$ the true
  $(n_1,n_2)$?  This is what an experiment uses, and it is much stricter: one
  gap in $B$ can shift every label behind it.
\end{enumerate}

\section{Which devices}

\begin{center}
\begin{tabular}{lrl}
\toprule
set & devices & used for \\
\midrule
training & \vNFit & fitting the U-Net weights \\
validation & \vNVal & choosing the epoch and the threshold $t$; choosing every
  post-processing setting \\
test & \vNTest & \textbf{only} reporting; nothing is chosen on it \\
\midrule
total pool & \vNPool & simulated once; split by device \\
\bottomrule
\end{tabular}
\end{center}

\noindent
The validation devices are \vValFrac\,\% of the \vNTrain{} training devices,
carved out with a fixed seed.  Because they chose the threshold, they are
\emph{not} held out; validation tables are reported but labelled as such.
The test devices are separate device IDs, and the split was checked to share
no device with the training set.

\paragraph{What was chosen on validation, never on test.}
\begin{itemize}
\item the U-Net threshold $t$, from \{\vThresholds\}, by F1@1 (stored in
  each checkpoint; \vThr{} for \vBud);
\item the dense no-ML detector: \vNDense{} settings, by state recall;
\item the line fit: \vNFitSet{} settings, by few-electron balanced accuracy
  (chosen: \emph{\vFitKnob}).
\end{itemize}

\paragraph{One checkpoint per budget.}  Some budgets were retrained with
another seed after the first training collapsed.  Every number here uses the
checkpoint the charge-state study used (for \vBud: \texttt{\vModel}), and
the line-level metrics were recomputed from it.

\section{Line-level metrics}

A line is one pixel wide, so pixel-by-pixel comparison fails twice:
\begin{itemize}
\item \textbf{pixel accuracy is useless}: predicting no line anywhere
  already scores \vNoLineAcc\,\%, because only \vMeTrue\,\% of the pixels are
  lines;
\item \textbf{strict pixel matching is too harsh}: a line drawn one pixel to
  the side scores zero.
\end{itemize}

\paragraph{Tolerant precision, recall and F1 at $\tau$ pixels.}  With $d_Y$
the distance from a pixel to the nearest true line pixel and $d_B$ the
distance to the nearest predicted one:
\begin{align}
\text{precision}@\tau &= \frac{\#\{p\in B: d_Y(p)\le\tau\}}{\#B}, &
\text{recall}@\tau &= \frac{\#\{p\in Y: d_B(p)\le\tau\}}{\#Y}, &
\text{F1}@\tau &= \frac{2PR}{P+R}.
\end{align}
$\tau=0$ is strict matching; $\tau=1$ is the headline (a line found within
one pixel).  Figure~\ref{fig:line} shows what is counted on \vDevice:
precision \vLPzero, recall \vLRzero{} at $\tau=0$, and precision \vLPone,
recall \vLRone, F1 \vLFone{} at $\tau=1$.

\begin{figure}[H]
\centering
\includegraphics[width=\textwidth]{figures/eval_line.pdf}
\caption{Line-level scoring on \vDevice, \vBud{} U-Net.  (a) What was
measured.  (b) U-Net probability.  (c, d) Each predicted and true line pixel
coloured by how it is counted, at $\tau=0$ and $\tau=1$.}
\label{fig:line}
\end{figure}

\paragraph{Reading F1@$\tau$ honestly.}  At $\tau$ the score credits a whole
disc of radius $\tau$ around each predicted pixel.  Two numbers keep this
honest:
\begin{itemize}
\item \textbf{coverage@$\tau$}: the fraction of the plane within $\tau$ of a
  \emph{measured} pixel (\vMeCov\,\% at $\tau=0$, \vMeCovOne\,\% at $\tau=1$
  for \vBud);
\item \textbf{predicted line fraction}: \vMePred\,\% of the pixels are
  predicted as line, against \vMeTrue\,\% true.  The U-Net over-draws, which
  raises recall and lowers precision.
\end{itemize}

\begin{table}[H]
\centering
\caption{Line-level metrics on the test devices, every U-Net of the sweep,
recomputed from the checkpoints the charge-state study used.  Means over
devices.}
\small
\resizebox{\textwidth}{!}{\input{line_table.tex}}
\end{table}

\paragraph{Why the line level is not enough.}  F1@1 rewards finding most of
every line, but the charge states need every line to be \emph{closed}.  A
map with F1@1 = \vMeFone{} (the mean for \vBud) can still merge two states
through one gap, and then every label behind that gap is wrong.  So the
model is also evaluated on the states themselves.

\section{Charge-state metrics}

All are computed per device and then averaged over devices.  Pixels within
1\,px of a true line are \emph{not scored}: their state is ambiguous at this
resolution, which matches the 1-pixel tolerance of F1@1.  On \vDevice,
\vDevScored\,\% of the pixels are scored.

\subsection*{Whole window}
\begin{itemize}
\item \textbf{pixel accuracy, relative}: apply the constant shift
  $(\Delta n_1,\Delta n_2)$ that makes the most pixels right, then take the
  fraction of scored pixels with the right $(n_1,n_2)$.  This tests the
  pattern, not the absolute count.
\item \textbf{state recall}: the fraction of the true states whose scored
  pixels are mostly ($>50\,\%$) labelled right.  \textbf{State precision}:
  the same from the side of the predicted regions.
\item \textbf{(1,1) IoU and centre error}: the overlap and centroid distance
  (mV) of the predicted and true $(1,1)$ regions.
\item \textbf{consistency}: the fraction of borders outside the labelling
  tree whose charge step agrees with the labels.  It uses no ground truth,
  so it can be computed in an experiment.
\end{itemize}

\subsection*{Few-electron region (the main evaluation)}
An experiment works near $(0,0)$.  Only the pixels whose \emph{true} state
has $n_1,n_2\le\vNmax$ are scored, with \emph{absolute} labels (no shift),
because the experiment needs the real occupation.
\begin{itemize}
\item \textbf{balanced accuracy} (headline): the accuracy per true state,
  averaged over the states present.  Every state counts once, so the large
  $(0,0)$ area cannot dominate.
\item \textbf{(1,1) landing} (headline): set the gates at the centre of the
  predicted $(1,1)$ region, its point farthest from the region's edge.  It
  scores 1 if the device is really in $(1,1)$ there, and 0 if not or if no
  $(1,1)$ was predicted.  This is the yes/no question of tuning.
\item \textbf{pixel accuracy}: reported, but dominated by $(0,0)$.
\end{itemize}

\begin{figure}[H]
\centering
\includegraphics[width=\textwidth]{figures/eval_states.pdf}
\caption{Charge-state scoring on \vDevice, \vBud{} U-Net with the line fit.
(a) Truth.  (b) Predicted labels.  (c) Scored pixels right (grey) or wrong
(red), unscored pixels white; the few-electron region is outlined.  (d) The
predicted $(1,1)$ region (outline) and the landing point (star).  (e) The
accuracy of each true state in the few-electron region; the balanced
accuracy is their mean (dashed).  On this device: relative accuracy
\vDevAccRel\,\%, state recall \vDevRec\,\%, few-electron pixel accuracy
\vDevAccFew\,\%, balanced \vDevBal\,\% over \vDevNSt{} states, landing:
\vDevLand.}
\label{fig:states}
\end{figure}

\section{References every number is read against}

A metric means little alone.  Every table carries three reference rows,
produced by the same scoring code:
\begin{itemize}
\item \textbf{all (0,0)} (the floor): label every pixel $(0,0)$.  It shows
  how much a metric gives away for free.  In the few-electron region it has
  balanced accuracy \vBaseBal\,\% and landing \vBaseLand\,\%.
\item \textbf{full CSD} (a check, not a result): label the true line map $Y$.
  Because $Y$ is derived from the true states, this only shows that the
  labelling code itself is right.
\item \textbf{dense, no ML} (the competitor): measure every pixel of the
  sensor image and find the lines without a network.  This is the
  conventional route the U-Net has to justify itself against.
\end{itemize}

\section{How results are reported, and how sure they are}

\begin{itemize}
\item \textbf{Per device, then averaged.}  Each device counts once, so one
  dense honeycomb does not dominate.
\item \textbf{Mean and median.}  The devices are very uneven (4 to 102 states
  per window), so both are given; a large gap between them means a few
  devices drive the mean.
\item \textbf{Uncertainty.}  With \vNTest{} test devices, a mean is uncertain
  by several points.  The 95\,\% intervals below come from resampling the
  test devices with replacement (10\,000 bootstrap draws).  Two methods whose
  intervals overlap strongly should not be called different.
\end{itemize}

\begin{table}[H]
\centering
\caption{Headline metrics with 95\,\% bootstrap intervals over the
\vNTest{} test devices: mean [low, high].  U-Net rows: \vBud.}
\small
\input{ci_table.tex}
\end{table}

\noindent
For \vBud, the line fit raises balanced accuracy from \vUBal\,\%
[\vUBalLo, \vUBalHi] to \vFBal\,\% [\vFBalLo, \vFBalHi].  The intervals do
not overlap, so the improvement is not a fluctuation of the test set.

\section{Limits of this evaluation}
\begin{itemize}
\item \textbf{Simulation only, noise-free.}  The truth is exact here and
  unavailable in an experiment; real noise would lower every method.
\item \textbf{Validation is not held out.}  It chose the threshold and every
  post-processing setting.  Only the test numbers are evidence.
\item \textbf{50 test devices.}  The intervals above are wide; small
  differences between budgets are within them.
\item \textbf{The corner is assumed to be $(0,0)$} for the absolute labels.
  This holds on every device of this pool, not necessarily in an
  experiment.
\item \textbf{One sweep} (4--8 rays $\times$ 40--60 points).
\end{itemize}

\end{document}
